\documentclass[11pt,a4paper]{article}
\usepackage[margin=1in]{geometry}
\usepackage{times}
\usepackage{latexsym}
\usepackage{graphicx}
\usepackage{booktabs}
\usepackage{amsmath}
\usepackage{amssymb}
\usepackage{multirow}
\usepackage{xcolor}
\usepackage{subcaption}
\usepackage[colorlinks=true,citecolor=blue,linkcolor=blue]{hyperref}
\usepackage{natbib}
\renewcommand{\bibsection}{\section*{References}}

\title{Structural Stagnation Probes: From Detection to Compute-Efficient\\Early Termination in LLM Agent Trajectories}

\author{Anonymous}

\begin{document}
\maketitle

%=============================================================================
\begin{abstract}
LLM agents frequently enter \emph{semantic stagnation}---executing diverse yet collectively ineffective action sequences during long-horizon tasks.
We present a framework for real-time stagnation monitoring that bridges detection with compute-efficient intervention.
Thirty lightweight structural features ($<$5\,ms per window) achieve 0.967 macro-F1 in zero-shot cross-agent transfer from SWE-Agent to OpenHands, outperforming the best LLM judge by 26.6\,pp in matched evaluation on identical OpenHands trajectories (0.967 vs.\ 0.701) and a zero-shot neural probe by 36.9\,pp.
Human gold validation (505 windows, 3 annotators, Fleiss $\kappa$=0.39) confirms that structural classifiers surpass majority-vote silver labels (GB 0.839 vs.\ MV 0.798, $p<0.001$).
On the 282 windows where labeling functions disagree (conf=1), the classifier maintains F1=0.784 versus MV agreement of 73.8\%, demonstrating learned decision boundaries beyond binary heuristics.
Three encoders (Qwen3-0.6B, BGE-M3, E5-large-v2) all degrade by 31--33\,pp under cross-agent transfer (mean gap to structural: 35.8$\pm$1.2\,pp), confirming an encoder-independent structural phenomenon.
Retrospective compute-waste analysis reveals that 52.4\% of agent steps occur after stagnation onset.
An early termination policy at $k$=3 consecutive stagnation windows saves 35.3\% of compute while preserving 95.5\% task success, with a Pareto frontier closely tracking the oracle.
Temperature-calibrated predictions achieve ECE$<$0.05, enabling reliable confidence-based deployment.
We additionally report a well-characterized negative result: three temporal label model variants fail to improve over majority vote, constraining future research on sequential weak supervision.
\end{abstract}

%=============================================================================
\section{Introduction}
\label{sec:intro}

LLM-based autonomous agents have demonstrated impressive capabilities on software engineering tasks~\cite{yang2024swe,wang2024opendevin}, yet they frequently exhibit a failure mode we term \emph{semantic stagnation}: the agent executes superficially diverse actions---different commands, different files, different strategies---that collectively make no progress toward the task goal.
Unlike simple command repetition, semantic stagnation evades na\"ive cycle detectors because individual actions appear novel.

Recent work has quantified this degradation at scale: SlopCodeBench~\cite{slopcodebench2026} reports structural erosion in 77\% of coding agent trajectories, with verbosity increasing in 75.5\%.
Process reward models (PRMs)~\cite{agentprm2026,prm_framework2025} address step-level evaluation but require outcome labels and RL pipelines, making them expensive training signals rather than lightweight monitors.
Self-reflection approaches~\cite{shinn2023reflexion,yao2023tree} rely on the stagnating model's own judgment---precisely the faculty that has failed.

We propose a framework that spans the full pipeline from detection to intervention:

\begin{enumerate}
\item \textbf{Structural detection dominates zero-shot transfer.}
We formulate stagnation detection as window-level binary classification over 30 structural features (command repetition rates, file edit counts, output variance, etc.).
A Gradient Boosting classifier trained on SWE-Agent with weak supervision achieves 0.967 cross-agent F1 on OpenHands \emph{with zero target-domain data}, outperforming the best LLM judge (Claude few-shot-CoT, 0.701) by 26.6\,pp in matched evaluation and a zero-shot neural probe (0.598) by 36.9\,pp.

\item \textbf{Gold-validated anti-circularity evidence.}
Since structural features overlap with labeling function (LF) inputs, high silver-label performance risks circularity.
We validate on 505 human-annotated gold windows: GB achieves 0.839 F1, exceeding MV silver labels (0.798) and the weakest human annotator (0.831).
On 282 conf=1 windows where LFs disagree, GB maintains 0.784 F1 versus MV agreement of 73.8\% ($p<0.01$), proving the classifier captures beyond-LF signals.

\item \textbf{Detection-to-intervention bridge.}
Retrospective analysis shows 52.4\% of steps occur post-stagnation onset across 5{,}092 trajectories.
Early termination at $k$=3 saves 35.3\% compute while preserving 95.5\% task success, with a Pareto curve closely tracking the oracle detector.

\item \textbf{Temporal label model: a well-motivated negative finding.}
Three temporal label model implementations (time-weighted voting, HMM-TLM, lag-aware HMM) all fail to improve over simple majority vote, despite systematic temporal dynamics in LF behavior.
\end{enumerate}

%=============================================================================
\section{Related Work}
\label{sec:related}

\paragraph{Agent Failure Detection.}
Cycle detectors such as CDCS+CDSA and SpecRA identify structural repetition but miss semantically diverse stagnation.
TrajAD detects anomalous trajectories but operates at the trajectory level rather than providing window-level monitoring.
SlopCodeBench~\cite{slopcodebench2026} benchmarks coding agent degradation, reporting structural erosion in 77\% of trajectories---motivating real-time detection.

\paragraph{Process Reward Models.}
AgentPRM~\cite{agentprm2026} trains step-level PRMs via TD/GAE to evaluate ``promise and progress,'' but requires outcome labels and an RL pipeline.
Related work includes general PRM frameworks~\cite{prm_framework2025}, ToolPRMBench~\cite{toolprmbench2026}, and trajectory-level reward modeling~\cite{traj_reward2026}.
Our approach differs fundamentally: we use \emph{window-level weak supervision} with no outcome labels, producing a monitoring tool rather than a training signal.

\paragraph{LLM Self-Reflection.}
Reflexion~\cite{shinn2023reflexion}, LATS~\cite{zhou2023lats}, and Tree of Thoughts~\cite{yao2023tree} leverage the agent's own reasoning for course correction.
These are intervention mechanisms that depend on the stagnating model's judgment.
Our structural detector is external and independent of the monitored agent.

\paragraph{Agent Monitoring.}
Recent work on monitoring LLM agents for sequentially contextual harm~\cite{seqharm2025} argues that sequence-level analysis is necessary---single-step evaluation is insufficient.
Reliable weak-to-strong monitoring~\cite{w2s_monitor2025} studies weak model supervision of strong agents for safety, complementary to our progress-focused monitoring.

\paragraph{Weak Supervision and Data Programming.}
The Snorkel paradigm~\cite{ratner2017snorkel} aggregates noisy labeling functions via generative models.
We extend this to sequential agent data, discovering that temporal LF dynamics do \emph{not} yield improved label models in the low-conflict regime (Section~\ref{sec:tlm}).

%=============================================================================
\section{Method}
\label{sec:method}

\subsection{Problem Formulation}

Given an agent trajectory $\tau = (a_1, o_1, a_2, o_2, \ldots, a_T, o_T)$ of actions and observations, we partition $\tau$ into overlapping windows of $k=5$ steps with stride 1.
Each window $w_i = (a_i, o_i, \ldots, a_{i+4}, o_{i+4})$ is classified as \textsc{stagnation} or \textsc{progress}.

\subsection{Labeling Functions}

We employ three labeling functions (LFs) in a weak supervision pipeline:

\begin{itemize}
\item \textbf{cmd\_rep}: Fires when the command repetition rate within a window exceeds a threshold. High precision (97.9\%) but agent-specific coverage (65.5\% fire rate on SWE-Agent vs.\ 5.4\% on OpenHands).
\item \textbf{file\_diff}: Fires when no file modifications are detected in the window. Cross-agent stable (fire rate 17.8\% on OpenHands).
\item \textbf{outcome\_backpropagation (OB)}: Propagates task outcome backward as distant supervision. Privileged---requires final outcome, unavailable at inference time.
\end{itemize}

Silver labels are assigned by majority vote (MV) over LF outputs.
A window receives confidence $\text{conf} \in \{1,2,3\}$ based on the number of agreeing LFs.
For deployment, a \textbf{2-LF configuration} (cmd\_rep + file\_diff, no OB) achieves 0.978 F1 at conf$\geq$2, with 100\% label agreement with the 3-LF configuration on the shared conf$\geq$2 subset.

\subsection{Structural Features}

We extract 30 structural features per window, organized into categories:

\begin{itemize}
\item \emph{Command patterns}: repetition rate, unique command ratio, command diversity entropy
\item \emph{File activity}: edit binary, file access pattern, number of distinct files touched
\item \emph{Output statistics}: length variance, similarity between adjacent outputs
\item \emph{Temporal signals}: step duration trend, action type distribution
\item \emph{Error indicators}: error message count
\end{itemize}

These features are computed in $<$5\,ms per window, enabling real-time monitoring.

\subsection{Classifiers}

We train Gradient Boosting (GB) and Random Forest (RF) classifiers on SWE-Agent silver labels (conf$\geq$2).
For neural baselines, we train GRU probes on frozen embeddings from three encoders: Qwen3-Embedding-0.6B, BGE-M3, and E5-large-v2 (all $>$0.90 in-domain).

\subsection{Early Termination Policy}

Given the detector output $\hat{y}_i \in \{0,1\}$ for window $w_i$, we define an early termination policy with threshold $k$: terminate execution when $k$ consecutive windows are predicted as stagnation:
\begin{equation}
\text{TERMINATE} \iff \exists\, i : \hat{y}_j = 1 \;\; \forall j \in [i, i+k-1]
\end{equation}

%=============================================================================
\section{Experiments}
\label{sec:experiments}

\subsection{Anti-Circularity Validation on Gold Labels}
\label{sec:gold}

\paragraph{Gold Set Construction.}
Three annotators labeled 505 windows sampled from SWE-Agent trajectories (stratified by confidence level), achieving Fleiss $\kappa$=0.39---consistent with definitional ambiguity in stagnation boundaries (comparable to fine-grained sentiment analysis, $\kappa$=0.40--0.55).

\paragraph{Conf=1 Error Decomposition.}
The strongest anti-circularity evidence comes from the 282 conf=1 windows where LFs disagree completely.
On these windows, MV defaults to predicting progress (agreement with gold: 73.8\%), while GB achieves F1=0.784 ($p<0.01$ via paired bootstrap).
Table~\ref{tab:conf1_decomp} decomposes outcomes into four quadrants.

\begin{table}[t]
\centering
\small
\begin{tabular}{lrrl}
\toprule
\textbf{Category} & \textbf{Count} & \textbf{\%} & \textbf{Pattern} \\
\midrule
GB$\checkmark$ MV$\times$ (beyond-LF) & 48 & 17.0 & Subtle stag., composite signal \\
Both $\checkmark$ & 176 & 62.4 & Clear progress \\
Both $\times$ & 28 & 9.9 & Ambiguous boundary \\
GB$\times$ MV$\checkmark$ & 30 & 10.6 & FP: strategy transition \\
\bottomrule
\end{tabular}
\caption{Error decomposition on 282 conf=1 gold windows. GB rescues 48/76 (63.2\%) of MV errors---the beyond-LF detection rate.}
\label{tab:conf1_decomp}
\end{table}

The beyond-LF detection rate of 63.2\% (48 out of 76 MV errors corrected by GB) demonstrates that the classifier has learned soft decision boundaries over composite structural signals that generalize beyond binary LF heuristics.
The 30 GB-wrong-MV-correct cases (10.6\%) are predominantly false positives during strategy transitions, consistent with the overall FP-dominant (65\%) error profile.

\paragraph{Human Ceiling Analysis.}
Table~\ref{tab:human_ceiling} contextualizes GB performance against human annotators.
The Spearman-Brown adjusted ceiling is 0.870; GB at 0.839 is only 3.1\,pp below, exceeding the weakest annotator (B3, 0.831).

\begin{table}[t]
\centering
\small
\begin{tabular}{lcc}
\toprule
 & \textbf{F1} & \textbf{Gap to Ceiling (pp)} \\
\midrule
Adjusted Human Ceiling & 0.870 & --- \\
Best Annotator (B2) & 0.881 & +1.1 \\
\textbf{GB (structural)} & \textbf{0.839} & \textbf{$-$3.1} \\
Mean Annotator & 0.853 & $-$1.7 \\
Weakest Annotator (B3) & 0.831 & $-$3.9 \\
MV Silver Labels & 0.798 & $-$7.2 \\
GRU Probe & 0.790 & $-$8.0 \\
\bottomrule
\end{tabular}
\caption{Human ceiling vs.\ model performance on gold set (N=505). GB approaches the annotation ceiling ($-$3.1\,pp) and exceeds the weakest individual annotator.}
\label{tab:human_ceiling}
\end{table}

\paragraph{Power Analysis.}
With $N$=505 paired bootstrap samples, we can detect $\geq$4\,pp differences at $\alpha$=0.05 with power=0.82.
Table~\ref{tab:power} reports power for key comparisons.

\begin{table}[t]
\centering
\small
\begin{tabular}{lccc}
\toprule
\textbf{Comparison} & \textbf{$\Delta$pp} & \textbf{Power} & \textbf{Sig.?} \\
\midrule
GB vs MV & 4.2 & 0.82 & Yes ($p$=0.0001) \\
GB vs Probe & 8.8 & 0.99 & Yes ($p<$0.0001) \\
GB vs RF & 2.9 & 0.68 & Yes ($p$=0.0003) \\
TWV-Probe vs GB & $-$2.3 & 0.52 & No ($p$=0.056) \\
\bottomrule
\end{tabular}
\caption{Power analysis for gold set comparisons (N=505, $\alpha$=0.05). Effects $<$3\,pp would require $\sim$1{,}900 samples.}
\label{tab:power}
\end{table}

%-----------------------------------------------------------------------------
\subsection{Matched LLM Comparison on OpenHands}
\label{sec:llm_oh}

Prior comparisons between structural classifiers and LLM judges suffered from evaluation scope confounding: structural models were evaluated on SWE-Agent while LLMs were tested on different data.
We eliminate this by running LLM-as-judge (multiple prompting strategies) on the \emph{same} OpenHands trajectories used for structural cross-agent evaluation.

\begin{table}[t]
\centering
\small
\begin{tabular}{lcccc}
\toprule
\textbf{Method} & \textbf{F1} & \textbf{F1\textsubscript{stag}} & \textbf{F1\textsubscript{prog}} & \textbf{Cost} \\
\midrule
\textbf{Structural GB (OH)} & \textbf{0.967} & \textbf{0.948} & \textbf{0.987} & \$0 \\
Claude few-shot-CoT & 0.701 & 0.624 & 0.778 & \$48.3 \\
Claude few-shot & 0.662 & 0.581 & 0.743 & \$42.1 \\
GPT-5.5 zero-shot & 0.638 & 0.589 & 0.687 & \$12.8 \\
Claude zero-shot & 0.585 & 0.497 & 0.673 & \$13.2 \\
\bottomrule
\end{tabular}
\caption{Matched comparison on OpenHands trajectories (N=2{,}000 stratified sample). Structural GB outperforms the best LLM variant by 26.6\,pp at zero cost.}
\label{tab:llm_oh}
\end{table}

Table~\ref{tab:llm_oh} shows structural GB outperforms the best LLM (Claude few-shot-CoT) by 26.6\,pp on identical data.
LLMs additionally suffer cross-domain degradation: the best LLM drops 9.8\,pp from SWE to OpenHands (0.799$\to$0.701), while structural features degrade only 2.8\,pp (0.995$\to$0.967), demonstrating superior cross-domain robustness.

\begin{table}[t]
\centering
\small
\begin{tabular}{lccc}
\toprule
\textbf{Variant} & \textbf{SWE F1} & \textbf{OH F1} & \textbf{$\Delta$pp} \\
\midrule
Claude zero-shot & 0.663 & 0.585 & $-$7.8 \\
Claude few-shot-CoT & 0.799 & 0.701 & $-$9.8 \\
GPT-5.5 zero-shot & 0.747 & 0.638 & $-$10.9 \\
\bottomrule
\end{tabular}
\caption{LLM transfer degradation from SWE-Agent to OpenHands. All variants degrade 8--11\,pp, versus 2.8\,pp for structural features.}
\label{tab:llm_transfer}
\end{table}

%-----------------------------------------------------------------------------
\subsection{Encoder-Independence of the Cross-Agent Gap}
\label{sec:encoder}

A potential alternative explanation for the 36.9\,pp zero-shot gap between structural features and the Qwen3-0.6B probe is encoder weakness.
We test two additional encoders: BGE-M3 (BERT-based, 1024-d) and E5-large-v2 (contrastive-trained, 1024-d).

\begin{table}[t]
\centering
\small
\begin{tabular}{lcccc}
\toprule
\textbf{Encoder} & \textbf{In-domain} & \textbf{OH Zero-shot} & \textbf{$\Delta$pp} & \textbf{Gap vs Struct.} \\
\midrule
Qwen3-0.6B & 0.906 & 0.598 & $-$30.7 & 36.9\,pp \\
BGE-M3 & 0.941 & 0.621 & $-$32.0 & 34.6\,pp \\
E5-large-v2 & 0.936 & 0.607 & $-$32.9 & 36.0\,pp \\
\midrule
Structural GB & 0.995$^\dagger$ & 0.967 & $-$2.8 & --- \\
\bottomrule
\end{tabular}
\caption{Cross-agent zero-shot transfer by encoder. All three encoders degrade 31--33\,pp; gap to structural is 35.8$\pm$1.2\,pp. $^\dagger$Circular eval caveat: structural features overlap with LF inputs.}
\label{tab:encoder}
\end{table}

Despite BGE-M3 achieving +3.6\,pp over Qwen3 in-domain, its zero-shot gap shrinks by only 2.3\,pp (Table~\ref{tab:encoder}).
The cross-agent gap persists across all three encoders (34.6--36.9\,pp, mean 35.8$\pm$1.2\,pp), confirming it reflects a fundamental signal-type difference---structural vs.\ semantic---rather than encoder capability.

%-----------------------------------------------------------------------------
\subsection{Compute Waste Analysis}
\label{sec:waste}

We retrospectively analyze 5{,}092 trajectories (1{,}683 SWE-Agent + 3{,}409 OpenHands) to quantify computation expended after the first detected stagnation onset.

\begin{table}[t]
\centering
\small
\begin{tabular}{lrccc}
\toprule
\textbf{Agent} & \textbf{Traj.} & \textbf{Steps Wasted} & \textbf{Tokens Wasted} & \textbf{Med. Tail} \\
\midrule
SWE-Agent & 1{,}683 & 45.9\% & 49.4\% & 8 win. \\
OpenHands & 3{,}409 & 57.0\% & 60.0\% & 12 win. \\
\midrule
Combined & 5{,}092 & 52.4\% & 56.3\% & 11 win. \\
\bottomrule
\end{tabular}
\caption{Compute waste by agent. 78.6\% of trajectories contain at least one stagnation tail. The 90th percentile tail length is 38 windows (190 steps).}
\label{tab:waste}
\end{table}

Table~\ref{tab:waste} reveals that over half of all computation (52.4\% of steps, 56.3\% of tokens) occurs after stagnation onset.
OpenHands shows higher waste (57.0\% steps) due to longer trajectories and more diverse tool-calling patterns that extend ineffective exploration.
The token waste exceeds step waste because stagnation-phase steps produce longer outputs (verbose error messages and repeated explanations).

%-----------------------------------------------------------------------------
\subsection{Early Termination Simulation}
\label{sec:termination}

We simulate early termination policies with consecutive-stagnation thresholds $k \in \{1,2,3,5,7\}$ using the structural GB detector on combined SWE+OH test trajectories.

\begin{table}[t]
\centering
\small
\begin{tabular}{cccccc}
\toprule
$k$ & \textbf{Prec.} & \textbf{Recall} & \textbf{Saved\%} & \textbf{Success\%} & \textbf{False Term.} \\
\midrule
1 & 0.823 & 0.912 & 46.2 & 89.3 & 10.7\% \\
2 & 0.871 & 0.845 & 40.1 & 93.2 & 6.2\% \\
\textbf{3} & \textbf{0.908} & \textbf{0.783} & \textbf{35.3} & \textbf{95.5} & \textbf{3.8\%} \\
5 & 0.942 & 0.671 & 26.4 & 97.2 & 1.9\% \\
7 & 0.961 & 0.582 & 18.8 & 98.1 & 1.1\% \\
\bottomrule
\end{tabular}
\caption{Early termination simulation. $k$=3 is the recommended operating point: 35\% savings with $<$5\% task success loss and 3.8\% false termination rate.}
\label{tab:termination}
\end{table}

The $k$=3 operating point (Table~\ref{tab:termination}) achieves the best cost-effectiveness: 35.3\% compute savings while preserving 95.5\% task success.
Stagnation episodes of 1--2 windows frequently self-resolve (explaining $k$=1's high false termination rate of 10.7\%), while $\geq$3 consecutive stagnation windows rarely recover.

\paragraph{Pareto Frontier.}
We compare structural GB against oracle (perfect stagnation detection) and random baselines across operating points.
The structural Pareto curve closely tracks the oracle (gap of only 4.0\,pp at 35\% savings) while dramatically outperforming random termination (+19.3\,pp at the same savings level).
The area under the Pareto curve is 0.891 for structural, 0.952 for oracle, and 0.623 for random---structural detection achieves 93.6\% of oracle quality.

\paragraph{Sensitivity to Detector Quality.}
At $k$=3, a random detector saves only 22.1\% compute with 81.4\% success preservation, while the oracle saves 42.8\% with 99.6\% preservation.
Structural GB (35.3\% / 95.5\%) sits much closer to oracle than to random, confirming that detector quality, not just the termination policy, drives the savings-success trade-off.

%-----------------------------------------------------------------------------
\subsection{Deployable 2-LF Configuration}
\label{sec:2lf}

For real-time deployment, outcome\_backpropagation (OB) is unavailable since task outcomes are unknown during execution.
The 2-LF configuration (cmd\_rep + file\_diff) maintains 0.978 F1 at conf$\geq$2 with 100\% label agreement with the 3-LF configuration.
However, coverage varies significantly across agents (Table~\ref{tab:2lf}).

\begin{table}[t]
\centering
\small
\begin{tabular}{lrrc}
\toprule
\textbf{Agent} & \textbf{Windows} & \textbf{2-LF conf$\geq$2} & \textbf{Coverage} \\
\midrule
SWE-Agent & 54{,}809 & 10{,}243 & 18.7\% \\
OpenHands & 150{,}652 & 5{,}482 & 3.6\% \\
\bottomrule
\end{tabular}
\caption{2-LF coverage analysis. OpenHands coverage is low due to cmd\_rep's low fire rate (5.4\%) on OH's diverse command patterns.}
\label{tab:2lf}
\end{table}

The 3.6\% OpenHands coverage reflects cmd\_rep's low fire rate (5.4\%) on OH's more diverse command patterns.
This motivates a hierarchical deployment: 2-LF conf$\geq$2 for high-confidence alerts, supplemented by the structural classifier's continuous probability output ($P(\text{stagnation})>0.9$) for full coverage.

%-----------------------------------------------------------------------------
\subsection{Calibration Analysis}
\label{sec:calibration}

Deployment-grade early termination requires well-calibrated confidence scores.
Table~\ref{tab:calibration} reports expected calibration error (ECE) before and after temperature scaling.

\begin{table}[t]
\centering
\small
\begin{tabular}{llcccc}
\toprule
\textbf{Model} & \textbf{Data} & \textbf{ECE\textsubscript{raw}} & \textbf{ECE\textsubscript{cal}} & \textbf{Brier\textsubscript{raw}} & \textbf{Brier\textsubscript{cal}} \\
\midrule
GB & Silver & 0.061 & 0.045 & 0.063 & 0.059 \\
GB & Gold & 0.083 & 0.048 & 0.142 & 0.128 \\
RF & Silver & 0.042 & 0.038 & 0.041 & 0.039 \\
\bottomrule
\end{tabular}
\caption{Calibration metrics. Temperature scaling ($T$=2.06 for GB-silver, $T$=1.85 for GB-gold) brings ECE below 0.05 for all configurations.}
\label{tab:calibration}
\end{table}

GB predictions are highly bimodal: 64.4\% of predictions exceed 0.9 and 21.7\% fall below 0.1, with only 13.9\% in the mid-range.
Raw predictions in the 0.4--0.7 range show overconfidence (45--63\% model confidence vs.\ 12--23\% actual accuracy).
Temperature scaling ($T$=2.06) reduces ECE from 0.061 to 0.045 on silver and from 0.083 to 0.048 on gold---both below the 0.05 deployment threshold, enabling reliable confidence-based termination decisions.

%=============================================================================
\section{Analysis}
\label{sec:analysis}

\subsection{Feature Importance Decomposition}
\label{sec:shap}

SHAP analysis decomposes structural feature contributions into three categories relative to LF overlap (Table~\ref{tab:shap}).

\begin{table}[t]
\centering
\small
\begin{tabular}{lcrrl}
\toprule
\textbf{Category} & \textbf{\#Feat.} & \textbf{SHAP} & \textbf{\%} & \textbf{Examples} \\
\midrule
LF-overlap & 2 & 0.270 & 40.1 & cmd\_rep\_rate, file\_edit \\
Partial-overlap & 8 & 0.175 & 26.0 & unique\_cmd, error\_count \\
Residual (beyond-LF) & 20 & 0.228 & 33.9 & output\_var, duration\_trend \\
\bottomrule
\end{tabular}
\caption{SHAP feature importance decomposition. Residual features (33.9\%) account for GB's beyond-LF signal on conf=1 gold windows.}
\label{tab:shap}
\end{table}

The top two features (cmd\_repetition\_rate: 0.142, file\_edit\_binary: 0.128) directly correspond to LF inputs, explaining the near-perfect silver-label performance (0.995).
However, 20 residual features contribute 33.9\% of total SHAP value---these features (output\_length\_variance, step\_duration\_trend, cmd\_diversity\_entropy, output\_similarity\_adjacent, action\_type\_distribution) have no direct LF mapping and provide the beyond-LF signal that enables GB's 0.784 F1 on conf=1 gold windows.
Permutation importance ranks correlate strongly with SHAP ($\rho$=0.91).

\subsection{Qualitative Trajectory Analysis}

We present three representative trajectories illustrating key phenomena:

\paragraph{Example 1: Clear stagnation.}
An 18-window trajectory where the agent enters grep$\to$edit$\to$test cycling at window 7, with all three LFs firing and GB detecting immediately.
This represents the ``easy'' case where structural repetition is directly observable.

\paragraph{Example 2: Semantic stagnation.}
A 25-window trajectory where stagnation begins at window 12, but cmd\_rep does \emph{not} fire---the agent tries different approaches (file reads, API calls, configuration changes) that are individually diverse but collectively ineffective.
GB detects at window 13 via composite signal from output\_length\_variance and file\_edit\_binary.
This is the core phenomenon motivating our approach: stagnation that evades repetition-based detectors.

\paragraph{Example 3: False positive.}
A 15-window trajectory where GB predicts stagnation at window 6 during a legitimate strategy transition (the agent pauses to read documentation before pivoting).
Recovery occurs at window 8.
This FP pattern---strategy transitions misidentified as stagnation---accounts for the majority (65\%) of GB errors.

\subsection{Stride Sensitivity}
\label{sec:stride}

Window stride affects training set size, bootstrap independence, and model performance.
Table~\ref{tab:stride} shows results across strides 1, 3, and 5.

\begin{table}[t]
\centering
\small
\begin{tabular}{ccccc}
\toprule
\textbf{Stride} & \textbf{Train} & \textbf{Probe F1} & \textbf{GB F1} & \textbf{Autocorr.} \\
\midrule
1 & 38{,}161 & 0.906 & 0.995 & 0.72 \\
3 & 12{,}720 & 0.893 & 0.994 & 0.34 \\
5 & 7{,}632 & 0.882 & 0.992 & 0.11 \\
\bottomrule
\end{tabular}
\caption{Stride sensitivity. Structural GB is robust ($<$0.4\,pp range); probe degrades 2.4\,pp due to reduced training data.}
\label{tab:stride}
\end{table}

Structural GB performance is virtually unchanged across strides ($<$0.4\,pp), while the probe degrades by 2.4\,pp at stride 5 primarily due to data reduction (7.6K vs.\ 38K windows).
Stride 1 (default) maximizes data; stride 5 gives nearly independent windows (autocorrelation 0.11) with wider bootstrap CIs (0.042 vs.\ 0.018).

%=============================================================================
\section{Temporal Label Model: A Negative Finding}
\label{sec:tlm}

LFs exhibit systematic temporal dynamics along agent trajectories: on SWE-Agent, cmd\_rep leads stagnation onset by 1.87 steps, file\_diff is concurrent ($-$1.2 steps), and OB lags by +6.7 steps.
On OpenHands, cmd\_rep shifts from leading to lagging (+3.32 steps).
This i.i.d.\ violation motivated temporal label models (TLMs).

Three implementations all failed to improve over MV:
\begin{enumerate}
\item \textbf{Time-Weighted Voting (TWV)}: Harmful for most temporal bandwidths $\sigma$. The only marginally positive $\sigma$=2 degrades downstream classifiers: in-domain $-$4.9\,pp, cross-agent $-$21.7\,pp.
\item \textbf{HMM-TLM}: Overpredicts the stagnation class.
\item \textbf{Lag-aware HMM}: Merely matches MV (0.894 vs.\ 0.894).
\end{enumerate}

The failure is attributable to the low-conflict regime: conf$\geq$2 windows (where $\geq$2 LFs agree) have 95.1\% consistency with gold labels.
TWV primarily changes low-confidence labels, but changes them in the wrong direction.
This constitutes a well-characterized negative result: temporal LF dynamics exist but cannot be profitably exploited for label quality improvement in the current 3-LF, low-conflict setting.

%=============================================================================
\section{Limitations}
\label{sec:limitations}

\paragraph{Agent Scope.}
We evaluate on two SE-domain agents (SWE-Agent, OpenHands). Generalization to non-SE agents (web browsing, scientific research) requires future validation.

\paragraph{Gold Set Size.}
$N$=505 is powered for $\geq$4\,pp effects (power=0.82) but underpowered for smaller effects. The TWV-Probe vs.\ GB comparison ($\Delta$=2.3\,pp, $p$=0.056) is borderline due to insufficient power, not absent effect.

\paragraph{2-LF Coverage on OpenHands.}
The deployable 2-LF configuration covers only 3.6\% of OpenHands windows at conf$\geq$2, due to cmd\_rep's low fire rate (5.4\%). Full coverage requires using the structural classifier's probability output directly, bypassing LF-based confidence.

\paragraph{Simulation vs.\ Online Deployment.}
Early termination results are based on offline simulation. Online deployment requires agent framework integration and may face distribution shift from the termination policy itself affecting agent behavior.

\paragraph{Annotator Agreement.}
Fleiss $\kappa$=0.39 (fair agreement) reflects genuine definitional ambiguity in stagnation boundaries, comparable to discourse relation labeling ($\kappa$=0.35--0.50). This bounds the achievable performance ceiling (estimated 0.870).

%=============================================================================
\section{Conclusion}
\label{sec:conclusion}

We present a framework for LLM agent stagnation monitoring spanning detection to compute-efficient intervention.
Lightweight structural features achieve 0.967 cross-agent F1 with zero target-domain data, validated against human annotations (GB 0.839 vs.\ human ceiling 0.870, gap 3.1\,pp).
The beyond-LF detection rate of 63.2\% on disagreement windows provides direct anti-circularity evidence.
In matched evaluation, structural features outperform the best LLM judge by 26.6\,pp at zero inference cost, with encoder-independence confirmed across three architectures (gap 35.8$\pm$1.2\,pp).
Early termination at $k$=3 saves 35.3\% compute while preserving 95.5\% task success, with a Pareto frontier at 93.6\% of oracle quality.
Temperature-calibrated predictions (ECE$<$0.05) and a 2-LF deployable configuration (no privileged outcome information) enable practical deployment.
The negative TLM finding---temporal LF dynamics exist but cannot improve labels in low-conflict regimes---constrains future temporal weak supervision research.

\bibliographystyle{plainnat}
\bibliography{references}

%=============================================================================
\appendix

\section{Extended Feature Importance}
\label{app:shap}

The top 10 features by mean absolute SHAP value:

\begin{table}[h]
\centering
\small
\begin{tabular}{lcl}
\toprule
\textbf{Feature} & \textbf{SHAP} & \textbf{Category} \\
\midrule
cmd\_repetition\_rate & 0.142 & LF-overlap \\
file\_edit\_binary & 0.128 & LF-overlap \\
unique\_cmd\_ratio & 0.089 & Partial-overlap \\
output\_length\_variance & 0.067 & Residual \\
step\_duration\_trend & 0.054 & Residual \\
error\_message\_count & 0.048 & Partial-overlap \\
cmd\_diversity\_entropy & 0.043 & Residual \\
file\_access\_pattern & 0.038 & Partial-overlap \\
output\_similarity\_adjacent & 0.035 & Residual \\
action\_type\_distribution & 0.031 & Residual \\
\bottomrule
\end{tabular}
\caption{Top 10 structural features by SHAP importance (N=1{,}824 test windows). Permutation importance ranking correlation: $\rho$=0.91.}
\label{tab:shap_full}
\end{table}

\section{Full Pareto Frontier Data}
\label{app:pareto}

\begin{table}[h]
\centering
\small
\begin{tabular}{ccc}
\toprule
\textbf{Compute Saved (\%)} & \textbf{Success Preserved (\%)} & \textbf{Detector} \\
\midrule
18.8 & 98.1 & Structural \\
26.4 & 97.2 & Structural \\
35.3 & 95.5 & Structural \\
40.1 & 93.2 & Structural \\
46.2 & 89.3 & Structural \\
\midrule
25.1 & 99.8 & Oracle \\
35.6 & 99.5 & Oracle \\
42.8 & 99.1 & Oracle \\
48.5 & 98.2 & Oracle \\
52.4 & 96.8 & Oracle \\
\midrule
10.0 & 92.1 & Random \\
20.0 & 84.5 & Random \\
30.0 & 76.2 & Random \\
40.0 & 68.1 & Random \\
50.0 & 59.7 & Random \\
\bottomrule
\end{tabular}
\caption{Full Pareto frontier data. Area under curve: structural 0.891, oracle 0.952, random 0.623.}
\label{tab:pareto_full}
\end{table}

\section{Stagnation Tail Length Distribution}
\label{app:tail}

For the 78.6\% of trajectories with at least one stagnation tail:
\begin{itemize}
\item 25th percentile: 4 windows (20 steps)
\item Median: 11 windows (55 steps)
\item 75th percentile: 22 windows (110 steps)
\item 90th percentile: 38 windows (190 steps)
\end{itemize}

Unresolved trajectories show substantially higher waste (58--66\%) than resolved ones (10--14\%), as resolved trajectories' stagnation episodes are typically brief.

\end{document}
